\begin{lemma}[Additivity of $[[\gamma]]$ over a resolution]
  Let $E$ be a metric space, $\gamma \in \Theta(E)$ (\noderef{Curve}), $\omega \in \mathcal{D}^1(E)$
  (\noderef{Form1}), $k \in \mathbb{N}$, and $s \colon \mathbb{N} \to \mathbb{R}$ monotone
  (non-decreasing) on $\set{0,\dots,k}$ with $s(0)=0$ and $s(k)=\length(\gamma)$. Then
  \[
    [[\gamma]](\omega) = \sum_{i=0}^{k-1} \int_{s(i)}^{s(i+1)} \omega.f(\gamma(t)) \cdot
      \mathrm{deriv}(\omega.\pi\circ\gamma)(t)\,dt
    .
  \]
  This is the current-additivity half of paper Lemma~2.1/eq.~(1129) (paper.tex line 1129: ``if
  $\gamma$ is a concatenation of curves $\gamma_1,\dots,\gamma_k$ then we may split the
  Riemann-Stieltjes integral \eqref{currentcurve} across the subintervals corresponding to the
  concatenated curves, so that $[[\gamma]]=[[\gamma_1]]+\dotsb+[[\gamma_k]]$ follows immediately''),
  specialized to a resolution into (possibly degenerate) edges rather than to an abstract
  concatenation, and phrased -- exactly as \noderef{MorreyResolutionAdditive} does for the measure
  $\mu_\gamma$ -- without first constructing the reparametrized restricted curves
  $\gamma|_{[s(i),s(i+1)]}$ as objects of $\Theta(E)$: the un-reparametrized integral
  $\int_{s(i)}^{s(i+1)} \omega.f(\gamma(t)) \cdot \mathrm{deriv}(\omega.\pi\circ\gamma)(t)\,dt$ used here
  equals $[[\gamma|_{[s(i),s(i+1)]}]](\omega)$ for the reparametrized restricted curve, since
  reparametrizing by the constant shift $u \mapsto s(i)+u$ changes the integration variable by
  translation (which changes neither the value of $\omega.f\circ\gamma$ at corresponding points
  nor, by the translation-invariance of the derivative, the value of
  $\mathrm{deriv}(\omega.\pi\circ\gamma)$ there) and correspondingly relabels the bounds of integration
  from $[0,s(i+1)-s(i)]$ to $[s(i),s(i+1)]$ -- exactly the substitution used explicitly in the
  proof of \noderef{CurrentReverse}.
\end{lemma}
\begin{proof}
  \emph{Integrability of the integrand on every subinterval.} By \noderef{CurveLipschitz}, $\gamma$
  is globally $1$-Lipschitz, and $\omega.\pi$ is $\mathrm{Lip}(\omega.\pi)$-Lipschitz
  (\noderef{Form1}), so $\omega.\pi\circ\gamma$ is globally $\mathrm{Lip}(\omega.\pi)$-Lipschitz,
  hence Lipschitz (so absolutely continuous, \texttt{LipschitzOnWith.absolutelyContinuousOnInterval})
  on every interval $[a,b]$; consequently $\mathrm{deriv}(\omega.\pi\circ\gamma)$ is interval-integrable on
  every $a..b$ (\texttt{AbsolutelyContinuousOnInterval.intervalIntegrable\_deriv}, \texttt{Preamble}).
  Also $\omega.f$ is Lipschitz (\noderef{Form1}), hence continuous, so $\omega.f\circ\gamma$ is
  continuous on $\mathbb{R}$, hence continuous (and bounded) on every compact $[a,b]$. The product
  of an interval-integrable function with a continuous function on $[a,b]$ is interval-integrable
  on $a..b$; so $t \mapsto \omega.f(\gamma(t))\cdot\mathrm{deriv}(\omega.\pi\circ\gamma)(t)$ is
  interval-integrable on $a..b$, for every $a,b \in \mathbb{R}$.

  \emph{Reduction to a general additivity statement.} We prove, by induction on $n$, that for every
  $n \in \mathbb{N}$ and every $t \colon \mathbb{N} \to \mathbb{R}$ monotone on $\set{0,\dots,n}$,
  \[
    \int_{t(0)}^{t(n)} \omega.f(\gamma(u)) \cdot \mathrm{deriv}(\omega.\pi\circ\gamma)(u)\,du
    = \sum_{i=0}^{n-1} \int_{t(i)}^{t(i+1)} \omega.f(\gamma(u)) \cdot \mathrm{deriv}(\omega.\pi\circ\gamma)(u)\,du
    . \tag{$\ast_n$}
  \]
  Given $(\ast_k)$ applied to $t=s$, substituting $s(0)=0$ and $s(k)=\length(\gamma)$ turns the
  left-hand side into $[[\gamma]](\omega)$ (\noderef{curveCurrent}) and yields exactly the claimed
  identity.

  \emph{Base case $n=0$.} Both sides concern the single point $t(0)=t(0)$: the left side is
  $\int_{t(0)}^{t(0)} (\cdots) = 0$ (an interval integral over a degenerate interval vanishes), and
  the right side is an empty sum, also $0$. So $(\ast_0)$ holds.

  \emph{Inductive step.} Assume $(\ast_m)$ holds for all monotone $t$ on $\set{0,\dots,m}$; let $t$
  be monotone on $\set{0,\dots,m+1}$. Restricting to the subset $\set{0,\dots,m}$ shows $t$ is also
  monotone on $\set{0,\dots,m}$, so $(\ast_m)$ applies to $t$ itself. By the integrability
  established above, $u \mapsto \omega.f(\gamma(u))\cdot\mathrm{deriv}(\omega.\pi\circ\gamma)(u)$ is
  interval-integrable on $t(0)..t(m)$ and on $t(m)..t(m+1)$, so the standard additivity of the
  interval integral over adjacent intervals
  (\texttt{intervalIntegral.integral\_add\_adjacent\_intervals}, \texttt{Preamble}) gives
  \[
    \int_{t(0)}^{t(m)} (\cdots)\,du + \int_{t(m)}^{t(m+1)} (\cdots)\,du
    = \int_{t(0)}^{t(m+1)} (\cdots)\,du
    .
  \]
  Combining this with $(\ast_m)$ (applied to the first term on the left) gives
  \[
    \int_{t(0)}^{t(m+1)} (\cdots)\,du
    = \sum_{i=0}^{m-1} \int_{t(i)}^{t(i+1)} (\cdots)\,du + \int_{t(m)}^{t(m+1)} (\cdots)\,du
    = \sum_{i=0}^{m} \int_{t(i)}^{t(i+1)} (\cdots)\,du
    ,
  \]
  which is exactly $(\ast_{m+1})$. This completes the induction, hence the proof.
\end{proof}
